\subsection*{Что такое машинное обучение?}
Перед тем как начинать какие"=то рассуждения, важно понимать, что такое машинное обучение и какие задачи ставит перед собой данная область.
\textit{Машинное обучение} (\texttt{ML}, МО) "--- наука, решающая задачи поиска закономерностей в ограниченных данных с помощью построения т.н. \textit{модели, или алгоритма обучения}. В рамках наших занятий мы будем рассматривать раздел \textit{обучения с учителем}.

Обучениие с учителем (\texttt{supervised learning}) представляет собой класс задач, в котором в рамках данных известна целевая переменная (\texttt{target}), которую алгоритм учится предсказывать. Пока непонятно, что такое <<учится предсказывать>>, но к этому мы вернёмся ниже. Важна сама суть: пусть нам необходимо предсказать цену квартиры, основываясь на таких критериях, как площадь, этаж, район и т.д. Сначала нам дают всю необходимую информацию вместе с ценой, чтобы мы смогли изучить, какова реальная цена в зависимости от набора признаков. Затем нам дают другой набор данных, с теми же признаками, но уже без цены, и наша задача "--- \textit{предсказать} эту стоимость. Целевая переменная (стоимость квартиры) выступает <<учителем>>, по которому мы изучаем влияние данных признаков на стоимость. Вышеописанный пример относится к задаче \textit{регрессии}.

\textit{Данные играют ключевую роль в машинном обучении.} Именно на основе данных модель выявляет закономерности и формирует свою предсказательную способность. Здесь важно понимать, что мы не можем гарантированно знать, как и откуда появились наши данные, поэтому не стоит рассматривать машинное обучение как таблетку от всех проблем "--- \textit{модели неизбежно допускают ошибки}, а насколько критичные, зачастую, зависит от качества данных (и ещё от нескольки критериев, которые слишком рано брать в голову). Поэтому важно не только построить алгоритм, но и тщательно понимать, с какими данными мы работаем, как они устроены и какие потенциальные сложности в себе несут. Об этом подробнее — в следующем разделе.